\chapter*{How to Use This Book}
\markboth{How to Use This Book}{How to Use This Book}
\addcontentsline{toc}{chapter}{How to Use This Book}

Every chapter uses the same coloured boxes, so that you always know what you are reading. Here is the legend.

\begin{objectives}
\textbf{Learning Objectives}: opens every chapter and lists exactly what you will be able to do at the end. Use it as a checklist when you revise.
\end{objectives}

\begin{pullquote}[--- a famous person]
Each chapter also starts with a short quotation to set the mood.
\end{pullquote}

\begin{keyidea}
\textbf{Key Idea}: the one thing to remember. If you remember nothing else from a section, remember this box.
\end{keyidea}

\begin{definition}[: Term]
\textbf{Definition}: the formal statement of a concept, written in the form you can copy into an exam answer.
\end{definition}

\begin{example}[: Worked example]
\textbf{Example}: a small problem solved step by step with real numbers. Cover the solution and try it yourself first.
\end{example}

\begin{explain}
\textbf{Reading the Figure}: explains what every axis, colour, arrow and symbol of the picture above means, and what you should conclude from it.
\end{explain}

\begin{warning}
\textbf{Common Mistake}: errors students often make, and how to avoid them.
\end{warning}

\begin{tryit}
\textbf{Try It Yourself}: short Python programs with their real outputs. Run them, change the numbers, and watch what happens.
\end{tryit}

\begin{exam}
\textbf{Exam Tip}: how a topic is usually asked and how to answer it for full marks.
\end{exam}

\procon{Name of the method}
{\item Each method ends with this strip
 \item Advantages: what the method is good at}
{\item Disadvantages: where it fails or what it costs}
{\item Applications: where it is used in practice}

\begin{modelans}[: example]
\textbf{Model Answer}: a complete or outline answer to one of the exam questions, showing the expected structure.
\end{modelans}

\begin{revision}
\textbf{Quick Revision}: a one-page table at the end of every chapter that summarises every topic in one line. Read it the night before the exam.
\end{revision}

\begin{summary}
\textbf{Chapter Summary}: the main points of the chapter in a short list, placed before the Exam Corner.
\end{summary}

\paragraph{The cartoon characters.} Faces and speech bubbles appear in many figures. A smiling face means ``this works well'', a worried face means ``this is a problem'', and an angry face means ``strict'' or ``blame''. Do not skip them: they carry the idea of the picture.

\paragraph{Colours in figures.} Blue, orange and green are used for the classes or clusters; red usually marks errors, boundaries or the thing to notice; purple marks the method's own steps.
